% 2.9 — brief en documento/investigacion/2-9-mcp.md
\subsection{MODEL CONTEXT PROTOCOL (MCP)}
\subsubsection{DEFINICIÓN Y ARQUITECTURA}

El \textit{Model Context Protocol} (MCP) es un protocolo abierto que permite la integración entre aplicaciones basadas en modelos de lenguaje de gran tamaño (\textit{large language models}, LLM) y fuentes de datos y herramientas externas, y provee una forma estandarizada de conectar esos modelos con el contexto que necesitan \parencite{MCPSpec2026}. \textcite{Anthropic2024MCP} lo liberó en noviembre de 2024 con una motivación concreta: aun los modelos más sofisticados quedan aislados de los datos, tras silos de información y sistemas heredados, y cada nueva fuente exigía una implementación a medida. \textcite{Hou2025} lo caracterizan como un estándar unificado, bidireccional y con descubrimiento dinámico, distinto de las integraciones manuales de API y de los \textit{plugins} unidireccionales por ser abierto e independiente de la plataforma. En diciembre de 2025 Anthropic donó el protocolo a la Agentic AI Foundation, un fondo dirigido bajo la Linux Foundation, para garantizar su neutralidad respecto de proveedores \parencite{MCPBlog2025AAIF}.

La arquitectura distingue tres participantes: los \textit{hosts}, aplicaciones LLM que inician las conexiones; los clientes, conectores dentro del host; y los servidores, servicios que proveen contexto y capacidades \parencite{MCPSpec2026}. El host crea y administra los clientes, controla sus permisos y aplica las políticas de seguridad y consentimiento; cada cliente mantiene una conexión aislada con un único servidor, y el servidor, que puede ser un proceso local o un servicio remoto, opera con responsabilidades acotadas y no debería poder leer la conversación completa ni observar a otros servidores \parencite{MCPSpec2025}. Los mensajes se codifican en JSON-RPC 2.0, un protocolo de llamada a procedimiento remoto (\textit{remote procedure call}, RPC) ligero, sin estado e independiente del transporte \parencite{JSONRPC2013}: una solicitud lleva método, parámetros e identificador, y la respuesta devuelve un resultado o un error con el mismo identificador. La especificación define dos transportes: \textit{stdio}, en el que el cliente lanza el servidor como subproceso, y \textit{Streamable HTTP}, en el que el servidor atiende a múltiples clientes por un único \textit{endpoint} HTTP \parencite{MCPSpec2025}. La revisión vigente hizo el protocolo sin estado: eliminó las sesiones y cada solicitud lleva su versión y sus capacidades \parencite{MCPSpec2026}.

Un servidor ofrece tres primitivas: los recursos (\textit{resources}), contexto y datos para el usuario o el modelo; los \textit{prompts}, mensajes con plantilla y flujos de trabajo; y las herramientas (\textit{tools}), funciones que el modelo ejecuta \parencite{MCPSpec2026}. \textcite{Ehtesham2025} las clasifican según quién las controla: las herramientas por el modelo, los recursos por la aplicación y los \textit{prompts} por el usuario. La herramienta es la primitiva central de este trabajo. Cada una declara nombre, descripción y un esquema de entrada en JSON Schema; el cliente las descubre con \texttt{tools/list}, las invoca con \texttt{tools/call} y recibe contenido no estructurado o estructurado, validable contra un esquema de salida \parencite{MCPSpec2025}. La especificación separa los errores de protocolo de los errores de ejecución, como una regla de negocio violada, y recomienda que siempre exista una persona con capacidad de denegar una invocación. El cuadro~\ref{tab:mcp-primitivas} resume las primitivas y el uso que se prevé darles.

\begin{table}[H]
  \centering
  \caption{Primitivas del servidor MCP y su uso previsto en el proyecto}\label{tab:mcp-primitivas}
  \begin{tabular}{@{}>{\RaggedRight\arraybackslash}p{0.14\textwidth}>{\RaggedRight\arraybackslash}p{0.16\textwidth}>{\RaggedRight\arraybackslash}p{0.28\textwidth}>{\RaggedRight\arraybackslash}p{0.36\textwidth}@{}}
    \toprule
    \textbf{Primitiva} & \textbf{Control} & \textbf{Definición} & \textbf{Uso previsto} \\
    \midrule
    Herramientas (\textit{tools}) & Modelo & Funciones que el modelo descubre e invoca, con esquema de entrada y salida & Consultas de lectura sobre el ERP (centros de costo, certificaciones, presupuesto) y sobre las predicciones del servicio FastAPI \\
    Recursos (\textit{resources}) & Aplicación & Contexto y datos que el host adjunta al modelo & No se contempla en el alcance; posible exposición futura de documentación del ERP \\
    \textit{Prompts} & Usuario & Mensajes con plantilla y flujos de trabajo & No se contempla; corresponden al asistente que queda como trabajo futuro \\
    \bottomrule
  \end{tabular}
  \fuente{Elaboración propia a partir de \textcite{MCPSpec2026} y \textcite{Ehtesham2025}.}
\end{table}

El protocolo se apoya en la investigación sobre uso de herramientas por modelos de lenguaje. \textcite{Schick2023} parten de que los modelos fallan en tareas básicas, como la aritmética o la consulta de hechos, y entrenan uno capaz de decidir qué API llamar, cuándo, con qué argumentos y cómo incorporar el resultado. \textcite{Qin2024} llevan el problema a más de 16\,000 API REST reales, y \textcite{Patil2024} identifican el modo de fallo característico: argumentos incorrectos y alucinación del uso de una API, que se mitiga cuando el modelo consulta documentación en tiempo de ejecución. Esos trabajos resuelven el lado del modelo; MCP estandariza el opuesto, cómo un servidor describe sus herramientas con esquemas explícitos y las ejecuta, de modo que cualquier host las descubra sin los adaptadores a medida que \textcite{Ehtesham2025} señalan como obstáculo entre marcos de agentes. Los mismos autores sitúan a MCP como el primer escalón, el acceso a herramientas, frente a los protocolos entre agentes, que quedan fuera del alcance de este trabajo.

\subsubsection{INTEGRACIÓN DE SISTEMAS EMPRESARIALES CON MODELOS DE LENGUAJE}

Un ERP acumula datos transaccionales estructurados, actualizados a diario y sujetos a permisos por usuario, y eso condiciona cómo ponerlos al alcance de un modelo de lenguaje. \textcite{Lewis2020} señalan que los modelos preentrenados almacenan conocimiento factual en sus parámetros, pero acceden a él y lo manipulan con precisión limitada, y que dar procedencia a sus respuestas y actualizar su conocimiento siguen siendo problemas abiertos. Su propuesta, la generación aumentada por recuperación (\textit{retrieval-augmented generation}, RAG), combina esa memoria paramétrica con un índice vectorial consultado por un recuperador neuronal y produce lenguaje más específico y factual. Ese patrón recupera pasajes por similitud. En un ERP, en cambio, la pregunta habitual (cuánto lleva certificado un centro de costo, qué probabilidad de sobrecosto estima el modelo) se responde con una consulta estructurada sobre la fuente vigente y con permisos acordes a quien pregunta, según el mecanismo de acceso que se defina en el Capítulo III\@. \textcite{Anthropic2024MCP} plantea que un desarrollador puede exponer sus datos mediante servidores MCP o construir las aplicaciones cliente que los consumen; este proyecto toma el primer rol, con lo que la lógica de negocio y el control de acceso permanecen en el ERP\@.

El servidor del proyecto corre en Docker como proceso independiente y lo consumen clientes remotos, por lo que el transporte que corresponde es Streamable HTTP\@. La especificación define su flujo de autorización para los transportes HTTP, mientras que para \textit{stdio} recomienda tomar las credenciales del entorno \parencite{MCPSpec2026}. En ese flujo el servidor MCP actúa como servidor de recursos OAuth 2.1, el cliente MCP como cliente OAuth 2.1 y el servidor de autorización emite los tokens. El servidor MCP debe implementar \textit{OAuth 2.0 Protected Resource Metadata} \parencite{RFC9728}: ante una solicitud sin credenciales responde 401 con una cabecera \texttt{WWW-Authenticate} que remite a su documento de metadatos, donde declara su servidor de autorización. El token viaja como \textit{bearer} en cada solicitud, nunca en la URL, y el servidor debe validar que fue emitido específicamente para él como audiencia (véase 2.8.1), conforme a \textit{Resource Indicators for OAuth 2.0} \parencite{RFC8707}; un token inválido recibe 401 y un alcance insuficiente, 403 \parencite{MCPSpec2026}.

De ahí se desprende una consecuencia de diseño. Si el servidor MCP consulta API ascendentes, actúa ante ellas como cliente OAuth con un token separado y no debe reenviar el que recibió del cliente MCP, para evitar el problema del delegado confundido \parencite{MCPSpec2026}. En el proyecto, el token que valida el servidor MCP lleva como audiencia su URL canónica y no puede viajar al ERP en NestJS ni al servicio en FastAPI; hace falta un segundo token hacia esos servicios, cuyo mecanismo se define en los Capítulos III y IV\@. Del lado del proveedor de identidad, la guía oficial de \textcite{Keycloak2026MCP} documenta el uso de Keycloak como servidor de autorización para MCP y declara que no soporta \textit{Resource Indicators}; la audiencia se fija con un \textit{client scope} y un \textit{Audience mapper} que incluye la URL del servidor MCP en el token, mientras el cliente MCP sigue enviando el parámetro \texttt{resource} porque la especificación lo exige aunque el servidor de autorización no lo procese. La verificación del token en sí se trata en la sección 2.8.

La autorización no agota la seguridad. \textcite{Hou2025} sistematizan las amenazas por fase del ciclo de vida del servidor y por tipo de atacante, entre ellas el robo de credenciales durante la invocación y los conflictos de nombres que llevan a ejecutar la herramienta equivocada o escalar privilegios. \textcite{Radosevich2025} demuestran que modelos de primera línea pueden ser inducidos a usar herramientas MCP para ejecutar código malicioso, tomar control remoto o robar credenciales, y proponen un auditor que genera casos adversarios a partir de lo que un servidor expone. La especificación reconoce que las herramientas representan ejecución de código arbitrario y establece que el servidor debe validar las entradas, controlar el acceso, limitar la tasa de invocaciones y sanear las salidas \parencite{MCPSpec2025}. En este proyecto eso se traduce en un servidor que expone únicamente consultas de lectura acotadas sobre el ERP y las predicciones, sin nómina ni datos personales, con esquemas estrictos y acceso delegado en los alcances de Keycloak. Con ello el servidor deja el ERP y las predicciones listos para un asistente conversacional futuro; construir ese asistente queda fuera del alcance.
